\section{Discapacidad visual}

\subsection{Concepto de discapacidad visual}
La discapacidad visual se refiere a cualquier condición que limite o reduzca la capacidad de una persona para percibir su entorno mediante el sentido de la vista. Esta limitación puede presentarse en distintos grados y manifestarse a través de diversas alteraciones en el sistema visual, afectando funciones como la agudeza visual, el campo visual, la percepción de colores o la sensibilidad al contraste.

De acuerdo con la Organización Mundial de la Salud (OMS), la discapacidad visual “se produce cuando una afección ocular afecta al sistema visual y a sus funciones relacionadas con la visión” \cite{who_blindness2026}. Esta definición incluye una amplia variedad de condiciones médicas que pueden afectar los ojos o las vías visuales del sistema nervioso central.

En términos generales, la discapacidad visual engloba todas las formas de pérdida de visión que no pueden ser corregidas completamente mediante lentes, cirugía u otros tratamientos médicos convencionales. Esto incluye desde una disminución leve en la capacidad visual hasta la pérdida total de la visión.

Las causas de la discapacidad visual pueden ser múltiples y abarcan tanto factores congénitos como adquiridos. Entre las causas más comunes se encuentran los errores refractivos no corregidos, cataratas, glaucoma, degeneración macular relacionada con la edad, retinopatía diabética, lesiones traumáticas y enfermedades neurológicas que afectan el procesamiento visual.

Además del aspecto médico, la discapacidad visual también debe entenderse desde una perspectiva social y funcional. No solo implica una alteración en el sistema visual, sino que también puede generar limitaciones significativas en la realización de actividades cotidianas como leer, desplazarse de forma segura, reconocer rostros, utilizar dispositivos tecnológicos o interactuar con el entorno urbano.

Por esta razón, actualmente se promueve un enfoque basado en la Clasificación Internacional del Funcionamiento, de la Discapacidad y de la Salud \cite{who_icf2001}, que considera tanto las condiciones médicas como los factores ambientales y sociales que influyen en la calidad de vida de las personas con discapacidad visual.

\subsubsection{Definición de discapacidad visual según organismos internacionales}
Diversos organismos internacionales han desarrollado definiciones y criterios para clasificar la discapacidad visual con el objetivo de facilitar su diagnóstico, tratamiento y seguimiento a nivel global. Entre las instituciones más relevantes que han establecido lineamientos en este ámbito se encuentran la Organización Mundial de la Salud (OMS), la International Agency for the Prevention of Blindness (IAPB) y diferentes organismos gubernamentales responsables de la salud pública.

Según la OMS, la discapacidad visual ocurre cuando una afección ocular afecta el funcionamiento del sistema visual de tal manera que la persona experimenta limitaciones significativas en su capacidad para realizar actividades que dependen de la visión \cite{who_blindness2026}. Estas afecciones pueden afectar diferentes componentes del sistema visual, incluyendo la retina, el nervio óptico, la córnea o el cerebro.

Desde una perspectiva epidemiológica, la discapacidad visual representa un problema importante de salud pública. La OMS estima que en el mundo existen más de 2.200 millones de personas con algún tipo de deterioro visual, de las cuales aproximadamente 1.000 millones presentan afecciones que podrían haberse prevenido o tratado adecuadamente \cite{who_blindness2026}. Estas cifras reflejan la magnitud global del problema y resaltan la importancia de implementar políticas de prevención, diagnóstico temprano y tratamiento oportuno.

En muchos países se utiliza el criterio establecido por la OMS para definir la denominada ceguera legal, la cual se considera cuando la agudeza visual en el mejor ojo es inferior a 0.05 (3/60) incluso después de aplicar corrección óptica \cite{once_bajavision}. Este criterio permite determinar el nivel de discapacidad visual que puede dar acceso a programas de apoyo, servicios de rehabilitación y beneficios sociales.

Asimismo, algunas legislaciones consideran que una persona presenta discapacidad visual cuando su agudeza visual es inferior a 0.3 (6/18) en el mejor ojo con corrección. Este enfoque reconoce que incluso niveles moderados de pérdida visual pueden generar dificultades significativas en la vida diaria \cite{once_bajavision}.

En síntesis, los organismos internacionales coinciden en que la discapacidad visual debe definirse no solo en términos médicos, sino también en función de su impacto en la funcionalidad y en la capacidad de la persona para interactuar con su entorno.

\subsubsection{Clasificación de la discapacidad visual}
La clasificación de la discapacidad visual permite categorizar el grado de deterioro visual de una persona con base en criterios médicos estandarizados. Esta clasificación se basa principalmente en dos parámetros fundamentales:
\begin{itemize}
    \item Agudeza visual (AV)
    \item Campo visual
\end{itemize}

La agudeza visual mide la capacidad de una persona para distinguir detalles finos a una distancia determinada, mientras que el campo visual representa el área total que puede observar una persona sin mover los ojos.

De acuerdo con los criterios establecidos por la OMS y la Clasificación Internacional del Funcionamiento, la discapacidad visual puede dividirse en diferentes niveles de severidad según la agudeza visual del mejor ojo con corrección óptica \cite{once_bajavision}:
\begin{itemize}
    \item \textbf{Leve:} Agudeza visual entre 0.3 y 0.5 ($\geq$6/18). En este nivel, la persona puede realizar muchas actividades cotidianas, aunque puede experimentar dificultades en tareas que requieren precisión visual, como leer letras pequeñas o reconocer detalles a distancia.
    \item \textbf{Moderada:} Agudeza visual entre 0.1 y 0.3 ($<$6/18). Las personas con este grado de discapacidad visual presentan limitaciones más evidentes en actividades como leer, conducir o identificar objetos a distancia.
    \item \textbf{Severa:} Agudeza visual entre 0.05 y 0.1 ($<$6/60). En esta etapa la persona tiene una visión muy reducida y requiere ayudas ópticas o tecnológicas para realizar muchas actividades cotidianas \cite{once_bajavision}.
    \item \textbf{Ceguera:} Agudeza visual inferior a 0.05 ($<$3/60) en el mejor ojo con corrección óptica \cite{once_bajavision}. En esta condición, la visión es extremadamente limitada o inexistente.
\end{itemize}

Además de la agudeza visual, el campo visual también puede utilizarse para determinar la presencia de discapacidad visual. Por ejemplo, una persona puede tener una agudeza visual relativamente buena, pero presentar un campo visual extremadamente reducido (visión en túnel), lo cual genera limitaciones similares a las de una persona con baja visión.

Estas clasificaciones permiten establecer criterios clínicos uniformes y facilitan la comparación de estadísticas epidemiológicas a nivel internacional.

\subsection{Tipos de discapacidad visual}
La discapacidad visual puede manifestarse de diferentes maneras dependiendo del tipo de alteración que afecte al sistema visual. En general, estas condiciones se clasifican en función del grado de pérdida visual y de la naturaleza de las limitaciones que generan en la vida cotidiana.
\subsubsection{Ceguera total}
La ceguera total representa la forma más severa de discapacidad visual. Se caracteriza por una pérdida prácticamente completa de la visión, lo que impide que la persona perciba formas, colores o movimientos en el entorno.

Desde un punto de vista clínico, la ceguera total se define generalmente como una agudeza visual inferior a 0.05 (3/60) en el mejor ojo, incluso después de aplicar corrección óptica \cite{once_bajavision}. También puede diagnosticarse cuando el campo visual se encuentra restringido a menos de 10 grados.

Las personas con ceguera total no pueden utilizar la visión como medio principal para orientarse o interactuar con su entorno. Por esta razón, dependen en gran medida de otros sentidos como el tacto y el oído para realizar actividades cotidianas.

Existen dos tipos principales de ceguera total:
\begin{itemize}
    \item Ceguera congénita:
Se presenta desde el nacimiento y puede ser causada por factores genéticos, infecciones durante el embarazo o malformaciones del sistema visual.
\item Ceguera adquirida:
Ocurre posteriormente debido a enfermedades o lesiones que afectan el sistema visual.
\end{itemize}

Las principales causas de ceguera incluyen enfermedades como:
\begin{itemize}
    \item Retinopatía diabética
    \item Glaucoma avanzado
    \item Infecciones oculares graves
    \item Traumatismos oculares
    \item daños neurológicos en el nervio óptico

\end{itemize} \cite{who_blindness2026}

Debido a la ausencia de información visual, las personas con ceguera total suelen utilizar herramientas de apoyo como bastones blancos, perros guía o tecnologías de asistencia para desplazarse de manera segura.

\subsubsection{Baja visión}
La baja visión se refiere a una condición en la que una persona presenta una reducción significativa de su capacidad visual, pero aún conserva cierto grado de visión útil. A diferencia de la ceguera total, las personas con baja visión pueden percibir algunos elementos visuales como luces, sombras, formas grandes o colores intensos.

Generalmente se considera baja visión cuando la agudeza visual se encuentra entre 0.05 y 0.3 en el mejor ojo, incluso después de utilizar corrección óptica \cite{once_bajavision}.

Las personas con baja visión pueden enfrentar dificultades para realizar tareas como:
\begin{itemize}
    \item Leer textos 
    \item Reconocer rostros
    \item Distinguir objetos a distancia
    \item Desplazarse en entornos con iluminación deficiente
\end{itemize}\cite{who_blindness2026}

Entre las causas más frecuentes de baja visión se encuentran:
\begin{itemize}
    \item Degeneración macular
    \item Gglaucoma
    \item Cataratas avanzadas
    \item Daño retiniano
    \item Enfermedades genéticas o degenerativas
\end{itemize}\cite{who_blindness2026}

Aunque estas condiciones limitan significativamente la capacidad visual, muchas personas con baja visión pueden mantener un nivel considerable de independencia mediante el uso de ayudas ópticas, tecnológicas y programas de rehabilitación visual.

\subsubsection{Deficiencias visuales parciales}
Además de la reducción en la agudeza visual, existen otras alteraciones que afectan la calidad de la visión sin necesariamente provocar una pérdida completa de la capacidad visual. Estas condiciones se conocen como deficiencias visuales parciales.
\begin{itemize}
    \item \textbf{Defectos del campo visual:} Algunas enfermedades o lesiones pueden provocar la pérdida parcial de la visión periférica o central. Por ejemplo, el glaucoma puede generar una condición conocida como visión en túnel, en la cual la persona pierde gradualmente su campo visual periférico \cite{who_blindness2026}.
    \item \textbf{Alteraciones en la percepción del color:} El daltonismo y otras alteraciones cromáticas afectan la capacidad para distinguir ciertos colores. Aunque estas condiciones no suelen provocar discapacidad severa, pueden generar dificultades en actividades que dependen del reconocimiento cromático \cite{once_bajavision}.
    \item  \textbf{Problemas de sensibilidad al contraste:} Algunas personas presentan dificultades para distinguir objetos cuando existe poca diferencia de brillo entre ellos. Esto puede complicar tareas como leer textos impresos o identificar escalones \cite{once_bajavision}.
\end{itemize}
Estas deficiencias visuales pueden generar limitaciones significativas en situaciones reales, especialmente en entornos urbanos complejos donde la información visual es esencial para la orientación.

\subsection{Impacto de la discapacidad visual en la movilidad}
La discapacidad visual afecta profundamente la capacidad de una persona para desplazarse de manera segura y autónoma. La visión es uno de los principales sentidos utilizados para orientarse en el espacio, evitar obstáculos y planificar rutas de desplazamiento.

Cuando esta capacidad se ve limitada, las personas pueden experimentar dificultades significativas para realizar actividades cotidianas como caminar por la ciudad, utilizar transporte público o desplazarse en entornos desconocidos.
\subsubsection{Barreras físicas en entornos urbanos}
Las ciudades modernas presentan múltiples obstáculos que dificultan la movilidad de las personas con discapacidad visual. Entre los problemas más comunes:
\begin{itemize}
    \item Aceras irregulares
    \item Escalones sin señalización
    \item Postes u objetos colocados en zonas de paso
    \item Falta de rampas accesibles
    \item Ausencia de pavimento táctil
\end{itemize}

Estas barreras pueden generar riesgos significativos de caídas o accidentes. Investigaciones recientes indican que muchas personas con discapacidad visual encuentran dificultades para desplazarse de manera segura e independiente debido a la presencia de estos obstáculos \cite{beyond_cane2022}.

\subsubsection{Dificultades de orientación y desplazamiento}
La orientación espacial depende en gran medida de la información visual disponible. Cuando esta información no está presente, las personas deben recurrir a otras estrategias como:
\begin{itemize}
    \item Referencias auditivas
    \item Memoria espacial
    \item Señales táctiles
\end{itemize}

Sin embargo, estas estrategias no siempre son suficientes en entornos urbanos complejos. Según estudios recientes, la falta de información visual puede provocar desorientación, incremento del tiempo de desplazamiento y mayor riesgo de accidentes \cite{beyond_cane2022}.
\subsubsection{Importancia de la accesibilidad urbana}
La accesibilidad urbana busca diseñar espacios públicos que puedan ser utilizados por todas las personas, incluyendo aquellas con discapacidad visual.
Ya existen ciertas medidas para personas con discapacidad como el pavimento táctil para guiar rutas peatonales, los semáforos con señales auditivas, señalización en braille o iluminación adecuada en calles y edificios.

Estas medidas permiten mejorar la independencia y seguridad de las personas con discapacidad visual al desplazarse en espacios públicos.

\subsection{Tecnologías de asistencia para personas con discapacidad visual}
Las tecnologías de asistencia han evolucionado significativamente en las últimas décadas. Actualmente existen múltiples herramientas diseñadas para ayudar a las personas con discapacidad visual a interactuar con su entorno y desplazarse de manera más segura.

\subsubsection{Métodos tradicionales de asistencia (bastón, perro guía)}
Los métodos tradicionales incluyen herramientas ampliamente utilizadas como el bastón blanco o el perro guía

El bastón permite detectar obstáculos cercanos mediante el contacto con el suelo, mientras que el perro guía puede ayudar a evitar obstáculos y seguir rutas conocidas.

Aunque estos métodos siguen siendo fundamentales, ofrecen información limitada sobre el entorno y requieren entrenamiento especializado \cite{beyond_cane2022}.

\subsubsection{Tecnologías digitales de asistencia}
Los avances en tecnología han permitido el desarrollo de soluciones más avanzadas como:
\begin{itemize}
    \item Sistemas de navegación basados en GPS
    \item Aplicaciones móviles de reconocimiento visual
    \item Dispositivos portátiles con cámaras
\end{itemize}

Estas tecnologías pueden identificar objetos, leer textos y proporcionar información auditiva al usuario sobre su entorno. \cite{wearable_ai_2025

\subsubsection{Dispositivos inteligentes para la movilidad asistida}
Los dispositivos inteligentes combinan sensores, cámaras y algoritmos de inteligencia artificial para proporcionar asistencia avanzada durante la movilidad.

Entre los dispositivos más innovadores se encuentran las gafas inteligentes con cámaras, los bastones electrónicos con sensores, dispositivos vestibles con retroalimentación háptica.

Investigaciones recientes han demostrado que los sistemas que combinan visión artificial, señales auditivas y vibraciones pueden mejorar significativamente la capacidad de navegación de las personas con discapacidad visual \cite{wearable_ai_2025}.